\editionPage{indices-defis}{Indications progressives}{Lire une aide, puis reprendre la recherche}{Les indices orientent la démarche ; les solutions restent dans la partie Corrigés.}
\begin{blocEdition}{indices-defis}
\editionEtape{D01 · Deux racines dans un intervalle}
\textbf{Indice 1.} Repérer le sommet et calculer la valeur du polynôme aux deux bornes.\par\smallskip
\textbf{Indice 2.} Une parabole tournée vers le haut, positive aux bornes et négative à son sommet intérieur, coupe l'axe deux fois entre les bornes.\par\smallskip
\editionEtape{D02 · Le seuil d’une moyenne}
\textbf{Indice 1.} Chercher le point fixe de la relation affine, puis soustraire ce nombre à chaque terme.\par\smallskip
\textbf{Indice 2.} Pour la moyenne, utiliser $S_{n+1}=S_n+u_n$. Au test de boucle, conserver simultanément le rang, la somme $S_n$ et le terme $u_n$.\par\smallskip
\editionEtape{D03 · Trois tangentes par un point}
\textbf{Indice 1.} Écrire la tangente au point d'abscisse t avant de remplacer x par 1.\par\smallskip
\textbf{Indice 2.} Pour compter les solutions, utiliser les trois intervalles de monotonie de g et traiter séparément les hauteurs -3 et -2.\par\smallskip
\editionEtape{D04 · La fenêtre la plus lumineuse}
\textbf{Indice 1.} Tracer au crayon uniquement le contour extérieur ; sa longueur est la somme de quatre morceaux.\par\smallskip
\textbf{Indice 2.} Éliminer h avec le périmètre, puis optimiser le trinôme en r.\par\smallskip
\editionEtape{D05 · Encadrer sans logarithme}
\textbf{Indice 1.} Étudier le signe de la dérivée de $e^x-1-x$.\par\smallskip
\textbf{Indice 2.} Pour la majoration, remplacer $x$ par $-t$ et inverser deux nombres positifs. Puis prendre $t=1/(2n)$.\par\smallskip
\end{blocEdition}

\editionPage{indices-defis-2}{Indications progressives}{Lire une aide, puis reprendre la recherche}{Les indices orientent la démarche ; les solutions restent dans la partie Corrigés.}
\begin{blocEdition}{indices-defis-2}
\editionEtape{D06 · Additionner sinus et cosinus}
\textbf{Indice 1.} Dans le repère, écrire $x^2+y^2=1$ et $x+y=k$.\par\smallskip
\textbf{Indice 2.} Calculer le produit $xy$ grâce au carré de la somme, puis utiliser un trinôme de somme $k$ et de produit $(k^2-1)/2$.\par\smallskip
\editionEtape{D07 · Un triangle sous deux contraintes}
\textbf{Indice 1.} Développer le produit scalaire en coordonnées avant de reconnaître un cercle.\par\smallskip
\textbf{Indice 2.} Sur ce cercle, remplacer $x^2+y^2$ par 4 dans les carrés des distances. Le second lieu devient une droite verticale.\par\smallskip
\editionEtape{D08 · Une famille de cercles tangents}
\textbf{Indice 1.} Compléter deux carrés, puis remplacer $y$ par $6-x$ pour obtenir un trinôme en $x$.\par\smallskip
\textbf{Indice 2.} Pour projeter le centre, se déplacer dans la direction du vecteur normal $(1;1)$ jusqu'à satisfaire $x+y=6$.\par\smallskip
\editionEtape{D09 · Choisir une règle de contrôle}
\textbf{Indice 1.} Commencer par séparer les groupes D et son complément ; multiplier ensuite les poids des deux tests à l'intérieur de chaque groupe.\par\smallskip
\textbf{Indice 2.} Pour U, le seul chemin accepté réunit deux tests négatifs. Pour comparer les règles, calculer une différence de deux fonctions affines du coût d'erreur.\par\smallskip
\editionEtape{D10 · Même moyenne, objectifs différents}
\textbf{Indice 1.} Calculer d'abord l'espérance et le deuxième moment avec les probabilités des issues, même si deux lots coïncident.\par\smallskip
\textbf{Indice 2.} Mettre la variance sous forme canonique ; pour le gain strictement positif, traiter les trois issues une à une.\par\smallskip
\end{blocEdition}

\editionPage{indices-defis-3}{Indications progressives}{Lire une aide, puis reprendre la recherche}{Les indices orientent la démarche ; les solutions restent dans la partie Corrigés.}
\begin{blocEdition}{indices-defis-3}
\editionEtape{D11 · La dispersion d’un groupe réuni}
\textbf{Indice 1.} Dans le calcul de la variance réunie, tous les écarts doivent être mesurés à la même moyenne totale.\par\smallskip
\textbf{Indice 2.} La moyenne des écarts à la moyenne d'un groupe est nulle ; c'est ce qui annule le terme croisé du développement.\par\smallskip
\editionEtape{D12 · Un algorithme qui encadre une racine}
\textbf{Indice 1.} Utiliser la stricte croissance pour relier le signe de $f(m)$ à la position de $m$ par rapport à la racine.\par\smallskip
\textbf{Indice 2.} La largeur est une suite géométrique de raison $1/2$ ; distinguer la largeur de l'intervalle et l'erreur du milieu.\par\smallskip
\editionEtape{D13 · Une garantie pour tous les x}
\textbf{Indice 1.} Le minimum peut être au sommet ou à une borne : commencer par situer le sommet selon m.\par\smallskip
\textbf{Indice 2.} Comparer le minimum à la valeur $P_m(1/2)$, puis chercher un paramètre qui atteint cette borne pour tous les $x$.\par\smallskip
\editionEtape{D14 · Qualité, coût et risque}
\textbf{Indice 1.} Calculer l'espérance avant de dériver ; la différence entre les deux gains est constante.\par\smallskip
\textbf{Indice 2.} Le signe de $g'$ se lit avec la croissance de $e^x$. À l'optimum, remplacer $e^{-\alpha}$ par $1/2$ simplifie les probabilités et la variance.\par\smallskip
\end{blocEdition}

